% Generated from locale/tr/content/proof-theory/normalization/reductions.tex; do not hand-edit.


\olfileid[tr]{pt}{nor}{red}

\olsection{İndirgeme Dönüştürmeleri}

\Log{N1i} sistemindeki !!a{proof} içinde bir kesmenin uzunluğu $1$ ise
bu kesme, tek bir !!{formula} oluşumundan oluşan bir kesittir ve söz
konusu !!{formula} hem !!a{introduction} çıkarımının sonucu hem de
!!a{elimination} çıkarımının büyük öncülüdür. Normalleştirme ispatında
bu tür kesmeler, alt-!!{proof}---yani böyle bir
!!{introduction}/!!{elimination} çiftiyle biten parça---yerine çıkarım
çiftinin ortadan kaldırıldığı yeni bir !!{proof} konarak ``indirgenir''.
Böyle bir değiştirme yeni
kesmelere yol açabilir. Bununla birlikte, yeni kesmelerin derecesinin
indirgenen kesmeninkinden düşük olduğuna kendimizi ikna edebiliriz.
Dolayısıyla yeni !!{proof}, ya daha düşük bir kesme uzunluğuna sahiptir
(çünkü uzunluğu $1$ olan bir maksimal kesme kaldırılmıştır) ya da
!!{proof} için kesme derecesi azalmıştır (söz konusu maksimal kesme tek
maksimal kesmeyse).

Örneğin, ilgili çıkarımlar \Intro{\lif} ve \Elim{\lif}, $!A \ident !B
\lif !C$ ise ve aşağıdaki gibi biten bir $\delta_1$ alt-!!{proof}
varsa:
\[
\AxiomC{$\Discharge{!B}{x}$}
\RightLabel{$\delta_2$}
\DeduceC{$!C$}
\DischargeRule{\Intro{\lif}}{x}
\UnaryInfC{$!B \lif !C$}
\AxiomC{}
\RightLabel{$\delta_3$}
\DeduceC{$!B$}
\RightLabel{\Elim{\lif}}
\BinaryInfC{$!C$}
\DisplayProof
\]
Bu kesmeyi kaldırmak için $\delta$ içindeki $\delta_1$ alt ispatını
aşağıdakiyle değiştiririz:
\[
\AxiomC{}
\RightLabel{$\delta_3$}
\DeduceC{$!B$}
\RightLabel{$\Subst{\delta_2}{\delta_3}{!B^x}$}
\DeduceC{$!C$}
\DisplayProof
\]
$\Subst{\delta_2}{\delta_3}{!B^x}$ ile, $\delta_2$ içindeki $x$
etiketli her $!B$ varsayımının yerine $\delta_3$ !!{proof} bütününün
konmasıyla elde edilen !!{proof} kastediyoruz. Bu !!{proof} artık $x$
etiketli, $!B$ biçiminde açık varsayımlara sahip değildir. Her zamanki
gibi, aynı etiketli varsayımlar olarak iki farklı !!{formula}
geçmediğini varsaydığımız için başka hiçbir varsayımı da $x$ ile
etiketlenmiş değildir. Böylece yeni !!{proof}, $\delta_1$ ile aynı açık
varsayımlara sahiptir; $\delta_2$ içinde varsayımları boşa çıkaran her
çıkarım, $\Subst{\delta_2}{\delta_3}{!B^x}$ içinde aynı varsayımları
boşa çıkarır ve $\delta_3$ içinde varsayımları boşa çıkaran her çıkarımın
olasılıkla birden çok kopyası oluşarak $\Subst{\delta_2}{\delta_3}{!B^x}$
içindeki karşılık gelen varsayımları boşa çıkarır.

Bütünüyle $\delta_2$ ya da $\delta_3$ içinde veya $\delta$ içindeki
$\delta_1$'in bütünüyle dışında kalan her kesme değişmeden kalır: Yine
aynı !!{formula} biçimlerini içerir ve aynı çıkarımlardan geçer; dolayısıyla
özellikle, $\delta$ içindeki karşılık gelen kesmeyle aynı dereceye ve
uzunluğa sahiptir. Uzunluğu $1$ olan bir maksimal kesmeyi kaldırdık;
böylece ya kesme uzunluğunu ya da ($\delta$ için kesme uzunluğu $1$ ise
ve bu tek maksimal kesmeyse) kesme derecesini azalttık.

Ancak iki şey olabilir:
\begin{enumerate}
  \item Bütün $!B^x$ varsayımlarını $\delta_3$ ile değiştirerek
  $\delta_3$ içindeki bütün kesmeleri çoğalttık. Bunlardan bazıları
  maksimal kesmeyse yeni !!{proof} için kesme uzunluğu, eski !!{proof}
  için kesme uzunluğundan daha büyük olabilir. Bunun gerçekleşmemesini
  garanti etmeliyiz. Bu nedenle indirgeme dönüştürmelerini yalnızca
  \emph{en sağdaki} maksimal kesmelere, yani $\delta$ boyunca indirgenen
  kesmenin sağında kalan hiçbir dalda maksimal kesme bulunmayacak
  nitelikteki kesmelere uygularız. $\delta_3$ içinde bir kesme olsaydı
  burada indirgenen kesmenin sağında kalırdı ve dolayısıyla en sağdaki
  bir kesme üzerinde çalışıyor olmazdık. Her zaman en sağdaki kesmeleri
  seçerek, !!{proof} için kesme uzunluğunu azalttığımızı ya da kesme
  derecesini bile düşürdüğümüzü garanti ederiz.
  \item $\delta_2$ içindeki bir $!B^x$ varsayımı !!a{elimination}
  çıkarımının büyük öncülüyse ve $\delta_3$'ün son çıkarımı
  \Elim{\lor}, \Elim{\lexists} ya da !!a{introduction} çıkarımıysa
  $\delta_3$'ü böyle bir varsayım noktasında $\delta_2$'ye aşılayarak
  yeni bir kesme üretmiş oluruz. $\delta_2$ içinde yalnızca varsayım
  olan şey artık uzunluğu $1$ olan bir kesme (!!a{introduction}
  çıkarımının sonucu ve !!a{elimination} çıkarımının büyük öncülü) ya
  da bir kesme kesitinin son !!{formula} oluşumudur (!!a{elimination}
  çıkarımının büyük öncülü ve \Elim{\lor} ya da \Elim{\lexists}
  sonucu). $!B^x$, \Elim{\lor} ya da \Elim{\lexists} çıkarımının küçük
  öncülüyse zaten bir kesitin ve belki bir kesme kesitinin ilk
  !!{formula} oluşumuydu. Yeni !!{proof} içinde bu kesit artık daha uzun
  olabilir. Ancak bu yeni kesmeyi ya da var olan kesme kesitini oluşturan
  !!{formula}, $!B$'dir ve $\depth{!B} < \depth{!B \lif !C}$ olur.
  Dolayısıyla kesmeler eklemiş ya da kesmelerin uzunluğunu artırmış
  olabilsek de bunların hiçbiri maksimal kesme değildir. Böylece ya
  kesme derecesini ya da aynı derecede kesmeler kalmışsa kesme
  uzunluğunu azaltmış oluruz.
\end{enumerate}

Uzunluğu $1$ olan bir kesmeyi indirgeme işlemine \emph{indirgeme
dönüştürmesi} (ya da \emph{dolambaç dönüştürmesi}) denir. Bazı kural
çiftlerine ilişkin kalıplar \olref{tab:reduction-conversions}'da
verilmiştir. (Kalanlar alıştırma olarak bırakılmıştır.)

\begin{table}
\begin{multline*}
\bottomAlignProof
\AxiomC{$\Discharge{!B}{x}$}
\RightLabel{$\delta_2$}
\shortDeduce
\DeduceC{$!C$}
\DischargeRule{\Intro{\lif}}{x}
\UnaryInfC{$!B \lif !C$}
\AxiomC{}
\RightLabel{$\delta_3$}
\shortDeduce
\DeduceC{$!B$}
\RightLabel{\Elim{\lif}}
\BinaryInfC{$!C$}
\DisplayProof
\quad \longrightarrow\quad
\shoveright{\bottomAlignProof
\AxiomC{}
\RightLabel{$\delta_3$}
\shortDeduce
\DeduceC{$!B$}
\RightLabel{$\Subst{\delta_2}{\delta_3}{!B^x}$}
\shortDeduce
\DeduceC{$!C$}
\DisplayProof}
\\[4ex]
\shoveleft{\bottomAlignProof
\AxiomC{}
\RightLabel{$\delta_2$}
\shortDeduce
\DeduceC{$!B$}
\AxiomC{}
\RightLabel{$\delta_3$}
\shortDeduce
\DeduceC{$!C$}
\RightLabel{\Intro{\land}}
\BinaryInfC{$!B \land !C$}
\RightLabel{\Elim{\land}[1]}
\UnaryInfC{$!B$}
\DisplayProof
\quad \longrightarrow\quad
\bottomAlignProof
\AxiomC{}
\RightLabel{$\delta_2$}
\shortDeduce
\DeduceC{$!B$}
\DisplayProof}\qquad
\shoveright{\bottomAlignProof
\AxiomC{}
\RightLabel{$\delta_2$}
\shortDeduce
\DeduceC{$!B$}
\AxiomC{}
\RightLabel{$\delta_3$}
\shortDeduce
\DeduceC{$!C$}
\RightLabel{\Intro{\land}}
\BinaryInfC{$!B \land !C$}
\RightLabel{\Elim{\land}[2]}
\UnaryInfC{$!C$}
\DisplayProof
\quad \longrightarrow\quad
\bottomAlignProof
\AxiomC{}
\RightLabel{$\delta_3$}
\shortDeduce
\DeduceC{$!C$}
\DisplayProof}
\\[4ex]
\bottomAlignProof
\AxiomC{}
\RightLabel{$\delta_2$}
\shortDeduce
\DeduceC{$!B$}
\RightLabel{\Intro{\lor}[1]}
\UnaryInfC{$!B \lor !C$}
\AxiomC{$\Discharge{!B}{x}$}
\RightLabel{$\delta_3$}
\shortDeduce
\DeduceC{$!D$}
\AxiomC{$\Discharge{!C}{y}$}
\RightLabel{$\delta_4$}
\shortDeduce
\DeduceC{$!D$}
\DischargeRule{\Elim{\lor}}{x\,y}
\TrinaryInfC{$!D$}
\DisplayProof
\quad \longrightarrow\quad
\shoveright{\bottomAlignProof
\AxiomC{}
\RightLabel{$\delta_2$}
\shortDeduce
\DeduceC{$!B$}
\RightLabel{$\Subst{\delta_3}{\delta_2}{!B^x}$}
\shortDeduce
\DeduceC{$!D$}
\DisplayProof}
\\[2ex]
\shoveleft{\bottomAlignProof
\AxiomC{}
\RightLabel{$\delta_2$}
\shortDeduce
\DeduceC{$!B(c)$}
\RightLabel{\Intro{\lforall}}
\UnaryInfC{$\lforall[x][!B(x)]$}
\RightLabel{\Elim{\lforall}}
\UnaryInfC{$!B(t)$}
\DisplayProof
\quad \longrightarrow\quad
\bottomAlignProof
\AxiomC{}
\RightLabel{$\Subst{\delta_2}{t}{c}$}
\shortDeduce
\DeduceC{$!B(t)$}
\DisplayProof}
\\[4ex]
\shoveright{\bottomAlignProof
\AxiomC{}
\RightLabel{$\delta_2$}
\shortDeduce
\DeduceC{$!B(t)$}
\RightLabel{\Intro{\lexists}}
\UnaryInfC{$\lexists[x][!B(x)]$}
\AxiomC{$\Discharge{!B(c)}{x}$}
\RightLabel{$\delta_3$}
\shortDeduce
\DeduceC{$!D$}
\DischargeRule{\Elim{\lexists}}{x}
\BinaryInfC{$!D$}
\DisplayProof
\quad \longrightarrow\quad
\bottomAlignProof
\AxiomC{}
\RightLabel{$\delta_2$}
\shortDeduce
\DeduceC{$!B(t)$}
\RightLabel{$\Subst{\Subst{\delta_3}{t}{c}}{\delta_2}{!B(t)^x}$}
\shortDeduce
\DeduceC{$!D$}
\DisplayProof}
\end{multline*}
\caption{\Log{N1i} için bazı indirgeme dönüştürmeleri.}
\ollabel{tab:reduction-conversions}
\end{table}

\begin{prob}
Ardından \Elim{\lnot} gelen \Intro{\lnot} için indirgeme
dönüştürmesini verin.
\end{prob}

Ele alınması gereken bir durum daha vardır. Kesme tanımı, kesmenin
uzunluğunun $1$ olduğu ve $!A$'nın hem \FalseInt{} sonucu hem de
!!a{elimination} çıkarımının büyük öncülü olduğu durumu içerir. Bu
durum, $!A$'nın !!a{introduction} kuralının sonucu olduğu duruma benzer
biçimde ele alınır. Örneğin aşağıdakini
\[
\bottomAlignProof
\AxiomC{}
\RightLabel{$\delta_2$}
\DeduceC{$\lfalse$}
\RightLabel{\FalseInt}
\UnaryInfC{$!B \lif !C$}
\AxiomC{}
\RightLabel{$\delta_3$}
\DeduceC{$!B$}
\RightLabel{\Elim{\lif}}
\BinaryInfC{$!C$}
\DisplayProof
\quad\text{şununla değiştirin:}\quad
\bottomAlignProof
\AxiomC{}
\RightLabel{$\delta_2$}
\DeduceC{$\lfalse$}
\RightLabel{\FalseInt}
\UnaryInfC{$!C$}
\DisplayProof
\]
$!A \ident (!B \lor !C)$ ya da $!A \ident \lexists[x][!B(x)]$
durumlarında biraz dikkatli olmalıyız. Örneğin aşağıdakini
\[
\bottomAlignProof
\AxiomC{}
\RightLabel{$\delta_2$}
\DeduceC{$\lfalse$}
\RightLabel{\FalseInt}
\UnaryInfC{$!B \lor !C$}
\AxiomC{$\Discharge{!B}{x}$}
\RightLabel{$\delta_3$}
\DeduceC{$!D$}
\AxiomC{$\Discharge{!C}{y}$}
\RightLabel{$\delta_4$}
\DeduceC{$!D$}
\DischargeRule{\Elim{\lor}}{x\,y}
\TrinaryInfC{$!D$}
\DisplayProof
\quad\text{şununla değiştirirsek}\quad
\bottomAlignProof
\AxiomC{}
\RightLabel{$\delta_2$}
\DeduceC{$\lfalse$}
\RightLabel{\FalseInt}
\UnaryInfC{$!D$}
\DisplayProof
\]
kesme !!{formula} olarak $!D$'yi taşıyan yeni bir kesme ortaya
çıkarabiliriz ve $\depth{!D} < \depth{!B \lor !C}$ olacağının garantisi
yoktur!{} Bu nedenle burada da \Intro{\lor}/\Elim{\lor} indirgemesinde
olduğu gibi ilerlemeli ve şunlardan biriyle değiştirmeliyiz:
\[
\bottomAlignProof
\AxiomC{}
\RightLabel{$\delta_2$}
\DeduceC{$\lfalse$}
\RightLabel{\FalseInt}
\UnaryInfC{$!B$}
\RightLabel{$\delta_3$}
\DeduceC{$!D$}
\DisplayProof
\quad\text{ya da}\quad
\bottomAlignProof
\AxiomC{}
\RightLabel{$\delta_2$}
\DeduceC{$\lfalse$}
\RightLabel{\FalseInt}
\UnaryInfC{$!C$}
\RightLabel{$\delta_4$}
\DeduceC{$!D$}
\DisplayProof
\]

